\section{结论}

本文研究了带开放腿张量网络的完整稠密物化问题。与传统近似收缩主要面向标量结果或局部查询不同，本文关注的是这样一类任务：最终结果必须以完整张量的形式被显式生成，因此方法设计不能通过改变输出形式来规避计算代价。基于这一设定，本文将优化重点放在内部收缩过程，并提出了 \method{} 框架。

围绕这一目标，本文将完整物化过程组织为一条统一流程：先对输出索引空间进行分块，再在网络交互图上构造划分树，并沿树递归维护保持输出边界显式存在的分隔状态；在状态传播过程中，通过自适应低秩压缩和隐式草图合并抑制内部接口在合并阶段的维度膨胀。因而，\method{} 近似的不是最终输出对象，而是生成该输出过程中必须经过的内部分隔状态。

实验结果表明，\method{} 在 Ring-8、Tree-8、Grid $3\times 3$ 与 Random-8 上能够形成清楚且可用的精度--时间折中：在保持较小相对误差的同时，相对直接 TNC 获得显著加速。进一步的阶段拆解与机制分析说明，端到端收益几乎完全来自收缩阶段；其中，隐式草图合并是更核心也更稳定的效率来源，而输出分块在较困难的网络上同样具有实质作用。强度扫描实验则进一步表明，近似强度可以作为控制误差与时间权衡的有效调节旋钮。

总体而言，本文说明了：当完整输出是不可改变的任务要求时，近似收缩的合理位置并不在输出端，而在内部收缩过程中传播的分隔状态。沿着这一思路，划分树递归、分隔状态表示与草图驱动压缩可以被组织成一个统一且有效的完整物化框架。未来工作可进一步面向更大规模和更复杂的任务，继续完善选择性近似、缓存复用和更丰富的草图策略，使这种面向完整物化的近似收缩方法在更广泛场景下发挥作用。
